\documentclass[12pt]{article}
\usepackage{notes-project}

\begin{document}

\begin{center}
  {\LARGE\bfseries Lecture 6: Limit of Functions and Proving Limits Systematically}\\[0.25em]
  {\large Mathematical Analysis}\\[0.15em]
\end{center} 

\section{Limit of Functions} 

\begin{dfn}[Existence of a Limit at a Point]{limit-of-functions}
  Under metric space \((A, d_A)\) and \((B, d_B)\), let \(f:A\to B\) be a function. Let $c$ be a limit point of \(A\). 

  If there exists a point \(\ell \in B\) such that \[\forall \varepsilon \in \apos{\RR}, \, \exists \delta \in \apos{\RR}, \, \forall x \in A, \, 0 < d_A(x,c) < \delta \implies d_B(f(x),\ell) < \varepsilon\] then we say that \(f\) has a limit at \(c\)
\end{dfn}

\begin{thm}[Uniqueness of a Function Limit]{function-limit-uniqueness}
  If a function $f$ has a limit at a point \(c\), then the limit is unique, ie
  \[\exists \ell : B, \, \forall \varepsilon \in \apos{\RR}, \, \exists \delta \in \apos{\RR}, \, \forall x \in A, \, 0 < d_A(x,c) < \delta \implies d_B(f(x),\ell) < \varepsilon\]
  \[\iff\]
  \[\exists! \ell : B, \, \forall \varepsilon \in \apos{\RR}, \, \exists \delta \in \apos{\RR}, \, \forall x \in A, \, 0 < d_A(x,c) < \delta \implies d_B(f(x),\ell) < \varepsilon\]
\end{thm}

\begin{proof}
  Assume that \[\exists \ell_1, \ell_2 : B, \, \forall \varepsilon \in \apos{\RR}, \, \exists \delta \in \apos{\RR}, \, \forall x \in A, \, 0 < d_A(x,c) < \delta \implies d_B(f(x),\ell_1) < \varepsilon, \, d_B(f(x),\ell_2) < \varepsilon\] where $\ell_1 \neq \ell_2$

  We know that 
  \[d_B(f(x),\ell_1) + d_B(f(x),\ell_2) < 2 \varepsilon\] 

  Then, according to \xref[triangle inequality]{lec03::axm:metric-triangle-inequality}, 
  \[d_B(f(x),\ell_1) + d_B(f(x),\ell_2) \leq d_B(\ell_1, \ell_2) < 2 \varepsilon\] 

  According to \xref[epsilon criteria]{lec04::cor:equality-challenge-replacement}, \[d_B(\ell_1, \ell_2) = 0\], ie \xref[$\ell_1 = \ell_2$]{lec03::axm:metric-separation}. Contradiction. So uniqueness is proven. 
\end{proof}

\begin{rem}[Board Comparison]{sequence-function-limit-challenges}
  For a sequence \(a:\NNOne\to X\) in a
  \xref[metric space]{lec03::def:metric-space} \((X,d)\) and a function
  \(f:A\to B\), compare the two challenges
  \begin{center}
    \renewcommand{\arraystretch}{1.25}
    \begin{tabular}{@{}lcc@{}}
       & \textbf{Sequence} & \textbf{Function} \\
      \hline
      Arbitrary tolerance
        & \(\forall\varepsilon\in\apos{\RR}\)
        & \(\forall\varepsilon\in\apos{\RR}\) \\
      Restrictor
        & \(\exists N\in\NNOne\)
        & \(\exists\delta\in\apos{\RR}\) \\
      Input restriction
        & \(n\geq N\)
        & \(0<d_A(x,c)<\delta\) \\
      Output requirement
        & \(d(a(n),\ell)<\varepsilon\)
        & \(d_B(f(x),\ell)<\varepsilon\) \\
    \end{tabular}
  \end{center}
  The output requirement must hold for every \(n\in\NNOne\) or
  \(x\in A\) satisfying its input restriction

  \begin{enumerate}[nosep]
    \item \(\varepsilon\) is arbitrary small
    \item \(\delta\) and \(N\) is used to restrict \(x\)
      st \(x\) is close enough to the limit
    \item \(\delta\) and \(N\) are all depends on \(\varepsilon\)
  \end{enumerate}
\end{rem}

\begin{conv}[The \(\varepsilon\)--\(\delta\) Challenge]{epsilon-delta-challenge}
  A value \(\ell\in B\) passes the
  \emph{\(\varepsilon\)--\(\delta\) challenge} for \(f\) at \(c\) if,
  for any proposed \(\varepsilon>0\), there exists a radius \(\delta>0\)
  such that for all \(x\in A\) with \(0<d_A(x,c)<\delta\),
  \(d_B(f(x),\ell)<\varepsilon\)
\end{conv}

\begin{thm}[Function Limit Challenge Replacement]{function-limit-challenge-replacement}
  Let \((A,d_A)\) and \((B,d_B)\) be
  \xref[metric spaces]{lec03::def:metric-space}, let \(f:A\to B\),
  let \(c\in A\) be a \xref[limit point]{lec03::def:limit-points} of \(A\),
  and fix \(\ell\in B\). Let \(g\) be a
  \xref[convergence challenge function]{lec04::def:convergence-challenge-function}.
  With balls taken in \((A,d_A)\),
  \[
    \begin{gathered}
      \bigl(\forall\varepsilon\in\apos{\RR},\,
      \exists\delta\in\apos{\RR},\,
      \forall x\in\ball{\delta}{c}\setminus\{c\},\quad
      d_B(f(x),\ell)<\varepsilon\bigr)\\
      \iff\\
      \bigl(\forall\varepsilon\in\apos{\RR},\,
      \exists\delta\in\apos{\RR},\,
      \forall x\in\ball{\delta}{c}\setminus\{c\},\quad
      d_B(f(x),\ell)<g(\varepsilon)\bigr)
    \end{gathered}
  \]
\end{thm}

\begin{proof}
  \todo[Prove both directions using positivity and arbitrarily small values of the challenge function]
\end{proof}

\section{Function Limit Operator}

\begin{dfn}[Limit of a Function at a Point]{function-limit-relation}
  Let \((X,d_X)\) and \((B,d_B)\) be
  \xref[metric spaces]{lec03::def:metric-space}, let \(A\subseteq X\)
  and \(f:A\to B\), and let \(c\in X\) be a
  \xref[limit point]{lec03::def:limit-points} of \(A\).
  We say that \(f\) \emph{has limit \(\ell\in B\) at \(c\)} if
  \[
    \begin{gathered}
      \forall\varepsilon\in\apos{\RR},\quad
      \exists\delta\in\apos{\RR},\quad\forall x\in A,\\
      0<d_X(x,c)<\delta\implies d_B(f(x),\ell)<\varepsilon
    \end{gathered}
  \]
\end{dfn}

\begin{dfn}[Limit Domain]{function-limit-domain}
  A point in the limit domain if and only if
  \begin{enumerate}[nosep]
    \item \(c\) is a \xref[limit point]{lec03::def:limit-points} of the domain \(A\)
    \item \(f\) \xref[has a limit]{def:limit-of-functions} at \(c\)
  \end{enumerate}
  \[
    \LimDom:(A\to B)\to\PowerSet(X)
  \]
  \[
    \LimDom(f)\defeq
    \left\{c\in X\;\middle|\;
      \begin{gathered}
        \forall r\in\apos{\RR},\quad
        \ball{r}{c}\cap\bigl(A\setminus\{c\}\bigr)\neq\emptyset\\
        {}\land{}\\
        \exists\ell\in B,\quad\forall\varepsilon\in\apos{\RR},\quad
        \exists\delta\in\apos{\RR},\\
        \forall x\in A,\quad
        0<d_X(x,c)<\delta\implies d_B(f(x),\ell)<\varepsilon
      \end{gathered}
    \right\}
  \]
  \par\smallskip\noindent{\footnotesize\texttt{LaTeX: \detokenize{\LimDom(f)}}}
\end{dfn}

\begin{prop}[Well-Definedness of the Function Limit]{function-limit-well-defined}
  With the setup of the \xref[limit relation]{def:function-limit-relation},
  for every \(f:A\to B\) and \(c\in\LimDom(f)\), there exists a unique
  \(\ell\in B\) such that \(f\)
  \xref[has limit \(\ell\) at \(c\)]{def:function-limit-relation}
\end{prop}

\begin{proof}
  \todo[Use membership in the limit domain for existence and the limit-point property with the metric axioms for uniqueness]
\end{proof}

\begin{dfn}[Function Limit Operator]{function-limit-op}
  \[
    \fnlim:\Pi f:(A\to B),\quad\LimDom(f)\to B
  \]
  % Characterization for f:A->B, c in LimDom(f), and ell in B.
  \[
    \fnlim(f)(c)=\ell
    \iff
    \begin{gathered}
      \forall\varepsilon\in\apos{\RR},\quad
      \exists\delta\in\apos{\RR},\quad\forall x\in A,\\
      0<d_X(x,c)<\delta\implies d_B(f(x),\ell)<\varepsilon
    \end{gathered}
  \]
  \par\smallskip\noindent{\footnotesize\texttt{LaTeX: \detokenize{\fnlim(f)(c)}}}
\end{dfn}

\begin{conv}[Function Limit Notation]{function-limit-notation}
  For the \xref[function limit operator]{def:function-limit-op}, write
  \[
    \fnlim_c(f)\defeq\fnlim(f)(c)
  \]
  \par\smallskip\noindent{\footnotesize\texttt{LaTeX: \detokenize{\fnlim_c(f)}}}
\end{conv}

\section{One-Sided Limits}

\begin{dfn}[Right-Hand Limit]{right-hand-limit}
  \todo[Specify the real domain and right-hand accumulation assumption, then give the limit condition for inputs approaching c from the right]
\end{dfn}

\begin{dfn}[Left-Hand Limit]{left-hand-limit}
  \todo[Specify the real domain and left-hand accumulation assumption, then give the limit condition for inputs approaching c from the left]
\end{dfn}

\begin{conv}[One-Sided Limit Notation]{one-sided-limit-notation}
  \todo[Introduce the notation for right-hand and left-hand limits]
\end{conv}

\begin{thm}[Uniqueness of One-Sided Limits]{one-sided-limit-uniqueness}
  \todo[State uniqueness under the corresponding one-sided accumulation assumption]
\end{thm}

\begin{proof}
  \todo[Adapt the uniqueness argument to the one-sided input restriction]
\end{proof}

\begin{thm}[Two-Sided and One-Sided Limits]{two-sided-one-sided-limits}
  \todo[State when a two-sided limit is equivalent to equal right-hand and left-hand limits, including the required domain assumptions]
\end{thm}

\begin{proof}
  \todo[Prove both directions]
\end{proof}

\begin{rem}[Error Bounds for One-Sided Limits]{one-sided-error-bounds}
  \todo[Adapt the shared error-bound method to the right-hand and left-hand input regions]
\end{rem}

\section{Limits at Infinity}

\begin{dfn}[Limit at Positive Infinity]{limit-at-positive-infinity}
  \todo[Assume the real domain is unbounded above and define a finite output limit as the input tends to positive infinity]
\end{dfn}

\begin{dfn}[Limit at Negative Infinity]{limit-at-negative-infinity}
  \todo[Assume the real domain is unbounded below and define a finite output limit as the input tends to negative infinity]
\end{dfn}

\begin{conv}[Limit Notation at Infinity]{limit-at-infinity-notation}
  \todo[Introduce the notation for limits as the input tends to positive or negative infinity]
\end{conv}

\begin{thm}[Uniqueness of Limits at Infinity]{limit-at-infinity-uniqueness}
  \todo[State uniqueness under the corresponding unbounded-domain assumption]
\end{thm}

\begin{proof}
  \todo[Adapt the uniqueness argument using the corresponding input tail]
\end{proof}

\begin{rem}[Error Bounds for Limits at Infinity]{infinity-error-bounds}
  \todo[Adapt the shared error-bound method using a real input threshold M(t)]
\end{rem}

\section{A Systematic Method for Proving Limits}

The shared error-bound method covers finite-value limits of sequences,
two-sided function limits, one-sided limits, and limits as the input tends
to infinity. The input restriction changes with the approach, while the
output tolerance is handled using
\xref[challenge functions]{lec04::def:convergence-challenge-function}.
Infinite output limits use arbitrarily large output thresholds instead

\begin{rem}[Shared Proof Method]{systematic-limit-method}
  \todo[Summarize the shared procedure: choose a candidate limit, derive an error bound, verify the challenge function, and supply the corresponding radius or threshold]
\end{rem}

\begin{rem}[Using the Ball Radius as the Parameter]{ball-radius-parameter}
  In the \xref[challenge-function method]{thm:function-limit-challenge-replacement},
  when the input-radius parameter itself supplies the witness \(\delta\),
  the following condition suffices
  \[
    \forall\varepsilon\in\apos{\RR},\quad
    \forall x\in\ball{\varepsilon}{c}\setminus\{c\},\quad
    d_B(f(x),\ell)\leq g(\varepsilon)
  \]
  Here \(\varepsilon\) is the input-radius parameter, and
  \(g(\varepsilon)\) is the output-error bound for a
  \xref[challenge function]{lec04::def:convergence-challenge-function} \(g\).
  A local restriction such as an input radius at most \(1\) can instead
  be handled by choosing \(\delta=\min(\varepsilon,1)\) and verifying
  the output bound on that smaller ball
\end{rem}

\end{document}
